\documentclass[11pt,a4paper]{article}
\usepackage[top=1.7cm,bottom=1.5cm,left=1.5cm,right=1.5cm,headsep=8pt,footskip=16pt]{geometry}
\usepackage{enumitem}
\usepackage{amssymb}
\usepackage{array}
\usepackage[table]{xcolor}
\usepackage{fancyhdr}
\usepackage{lastpage}
\usepackage{multicol}
\usepackage{polyglossia}
\setdefaultlanguage[calendar=gregorian, numerals=maghrib]{arabic}
\setotherlanguage{english}
\IfFontExistsTF{Amiri}{%
  \newfontfamily\arabicfont[Script=Arabic, Scale=1.08]{Amiri}
}{%
  \newfontfamily\arabicfont[Script=Arabic, Scale=1.08, Path=../fonts/, Extension=.ttf,
    UprightFont=*-Regular, BoldFont=*-Bold, ItalicFont=*-Italic,
    BoldItalicFont=*-BoldItalic]{Amiri}
}
\IfFontExistsTF{AlBattar}{%
  \newfontfamily\khatt[Script=Arabic, Scale=1.15]{AlBattar}
}{%
  \newfontfamily\khatt[Script=Arabic, Scale=1.15, Path=../fonts/,
    BoldFont=ae_AlBattar.ttf]{ae_AlBattar.ttf}
}
\renewcommand{\headrulewidth}{0.8pt}
\renewcommand{\footrulewidth}{0.8pt}
\fancyhf{}
\rhead{\khatt المدرسة العليا للأساتذة ببوزريعة}
\lhead{\khatt ملحقة بومرداس}
\lfoot{\khatt الأستاذ: د. بن عبد الرزاق يوسف}
\cfoot{\thepage/\pageref{LastPage}}
\pagestyle{fancy}
\parindent 0pt
\setlength{\columnsep}{0.6cm}
\setlength{\columnseprule}{0.4pt}
\setlist[enumerate,1]{label=\textbf{\arabic*.}, leftmargin=*, itemsep=5pt, topsep=2pt}
\begin{document}

\linespread{1.08}\selectfont
\begin{center}
{\LARGE\khatt\textbf{الاختبار الثالث في الذكاء الاصطناعي}}\\[4pt]
{\normalsize وضعيات إشكالية ــ أسئلة متعددة الاختيارات ــ الفصول من 1 إلى 6}
\end{center}
\vspace{-2pt}\rule{\textwidth}{0.6pt}\\[2pt]
\textbf{التعليمة:} يتضمن الاختبار خمس وضعيات إشكالية. اقرأ كل وضعية بعناية، ثم اختر التحليل أو التصرف الأنسب
من بين الاختيارات الخمسة؛ لكل وضعية إجابة صحيحة واحدة فقط.
اكتب \textbf{حرف} الإجابة الصحيحة في جدول الإجابات في آخر الاختبار.
لكل وضعية 4 نقاط (المجموع 20 نقطة).
\vspace{4pt}

\noindent
{\large\khatt\textbf{الوضعية 1}}\hfill{\small(الفصل 1)}\par\nopagebreak\vspace{3pt}
\setlength{\fboxsep}{7pt}\colorbox{gray!10}{\parbox{\dimexpr\linewidth-14pt}{سأل طالب ماستر أداة ذكاء اصطناعي عن مراجع حديثة في البلاغة الجديدة، فاقترحت خمسة عناوين بمؤلفيها ودور نشرها وسنواتها. وجد منها ثلاثة كما ذُكرت، ووجد الرابع كتابًا حقيقيًا لكن لمؤلف آخر، ولم يجد الخامس إطلاقًا في الفهارس. ويستعد لتسليم قائمة مراجعه لمشرفه بعد يومين. أي تفسير لما حدث، وما يترتب عليه، هو الأدق؟}}\par\nopagebreak
\begin{itemize}[leftmargin=1.0cm, labelsep=6pt, itemsep=3pt, topsep=4pt, parsep=0pt]
  \item[\textbf{أ)}] الأداة استرجعت المراجع من قاعدة بيانات فهرسية قديمة، فوقع الخلط في المؤلف واختفى الكتاب الخامس من الطبعات الحديثة؛ ولذلك يكفي البحث عن العنوانين المشكلين في فهارس أقدم لاستكمال توثيقهما.
  \item[\textbf{ب)}] الأداة أدرجت عمدًا مراجع غير دقيقة لتختبر قدرة الطالب على التحقق، لأن الأدوات الحديثة مصممة تعليميًا لتنمية التفكير النقدي، فتقدّم أحيانًا معلومات خاطئة عن قصد في المجال الأكاديمي.
  \item[\textbf{ج)}] الخطأ في المؤلف سهو بسيط لا يؤثر في قيمة المرجع، والكتاب الخامس موجود على الأرجح بلغة أخرى لأن الأداة تترجم العناوين، فيكفي البحث عنه بالإنجليزية ثم توثيقه بالعنوان العربي الذي اقترحته الأداة، وهو ما يحدث كثيرًا مع الكتب الأجنبية.
  \item[\textbf{د)}] الأداة تولّد نصًا مرجَّحًا لغويًا، ولذلك يُتحقق فقط من المراجع التي لا يُعثر على عناوينها، أما المراجع التي تظهر عناوينها في الفهارس فتُعتمد بكل بياناتها كما اقترحتها الأداة دون مراجعة إضافية، ما دام العنوان نفسه صحيحًا.
  \item[\textbf{هـ)}] الأداة تولّد نصًا مرجَّحًا لغويًا ولا تسترجع سجلات فهرسية، فقد تجمع عنوانًا حقيقيًا مع مؤلف غير صاحبه، أو تركّب عنوانًا لا وجود له؛ ولذلك يُتحقق من كل عنصر في المرجع، لا من وجود العنوان وحده.
\end{itemize}
\par\vspace{6pt}
\noindent
{\large\khatt\textbf{الوضعية 2}}\hfill{\small(الفصل 3)}\par\nopagebreak\vspace{3pt}
\setlength{\fboxsep}{7pt}\colorbox{gray!10}{\parbox{\dimexpr\linewidth-14pt}{يريد أستاذ عشرين سؤال إعراب بأسلوب فروضه: جملة من ثماني كلمات تقريبًا، تحت كلمة واحدة منها خط، والمطلوب إعرابها، ثم التصحيح في سطر واحد. جرّب طلبًا وصفيًا مفصلًا، فجاءت الأسئلة متفاوتة الطول والشكل وطريقة التصحيح. ويحتاج إلى أسئلة جاهزة للطباعة دون تعديل يدوي طويل. ما التعديل الأنجع لضبط الشكل المطلوب؟}}\par\nopagebreak
\begin{itemize}[leftmargin=1.0cm, labelsep=6pt, itemsep=3pt, topsep=4pt, parsep=0pt]
  \item[\textbf{أ)}] يضيف إلى الطلب ثلاثة أسئلة من فروضه مع تصحيحها نماذجَ، ويطلب أسئلة جديدة على منوالها في البنية والطول وشكل التصحيح مع تحديد الأبواب النحوية، لأن الأمثلة تنقل الشكل بدقة يصعب بلوغها بالوصف وحده.
  \item[\textbf{ب)}] يطلب مئة سؤال بالوصف السابق نفسه ويختار منها يدويًا عشرين سؤالًا مطابقًا، لأن زيادة عدد الأسئلة المطلوبة تجعل الأداة أكثر التزامًا بالشكل الموصوف، دون حاجة إلى تقديم أي نماذج من فروضه، وهو أسلوب شائع عند المستخدمين.
  \item[\textbf{ج)}] يضيف مثالًا واحدًا يختلف عمدًا عن أسلوبه المعتاد ليحفّز الأداة على التنوع، ويطلب عشرين سؤالًا على منواله، لأن نموذجًا واحدًا مختلفًا يكفي لتعليم الأداة الشكل ويمنعها من تكرار القالب نفسه.
  \item[\textbf{د)}] يطيل الوصف بشروط أكثر تفصيلًا عن عدد الكلمات وموضع الكلمة المسطرة، ويطلب من الأداة أن تكون دقيقة جدًا، لأن الأمثلة تقيّد إبداع الأداة فتنسخها حرفيًا، بينما يترك الوصف المفصل لها حرية أوسع، وهذا أنسب للمهام الدقيقة.
  \item[\textbf{هـ)}] يضيف ثلاثة أسئلة من فروضه نماذجَ، ويطلب من الأداة إعادة صياغتها بكلمات مختلفة، لأن إعادة صياغة النماذج نفسها تضمن تطابق الشكل، وتنتج عشرين سؤالًا جديدًا مختلفًا في المحتوى والأبواب النحوية.
\end{itemize}
\par\vspace{6pt}
\noindent
{\large\khatt\textbf{الوضعية 3}}\hfill{\small(الفصل 4)}\par\nopagebreak\vspace{3pt}
\setlength{\fboxsep}{7pt}\colorbox{gray!10}{\parbox{\dimexpr\linewidth-14pt}{طلبت أستاذة من أداة توليد صور رسم «عالِم في مختبره» لعشرة ملصقات عن تاريخ العلوم عند العرب والمسلمين، فجاءت الصور كلها لرجال بملامح ولباس أوروبي معاصر وأجهزة مخبرية حديثة. وتريد تعليق الملصقات في بهو المؤسسة الأسبوع القادم. وهي تحرص على أن تعكس الملصقات حقيقة تاريخية وصورة متوازنة. ما التحليل الأدق والتصرف الأنسب؟}}\par\nopagebreak
\begin{itemize}[leftmargin=1.0cm, labelsep=6pt, itemsep=3pt, topsep=4pt, parsep=0pt]
  \item[\textbf{أ)}] النتيجة تعكس انحيازًا في بيانات التدريب، فتكتب في الطلب «ارسم صورة غير منحازة» دون تفاصيل أخرى، لأن الأدوات الحديثة تفهم مفهوم الانحياز وتصححه وحدها، فلا حاجة إلى مراجعة الصور بعد ذلك.
  \item[\textbf{ب)}] النتيجة تعكس انحيازًا في بيانات التدريب التي يغلب فيها هذا التمثيل، فتعدّل الطلب بتحديد السياق التاريخي والتنوع المطلوب صراحة، ثم تراجع الصور للتأكد من دقتها التاريخية وتوازن تمثيلها قبل الطباعة.
  \item[\textbf{ج)}] النتيجة دقيقة لأن الأداة ترسم الصورة الأكثر شيوعًا للعالِم، والشيوع يعكس الواقع، فتعتمد الصور كما هي، وتكتفي بإضافة أسماء علماء عرب ومسلمين تحتها لتحقيق هدف الملصقات.
  \item[\textbf{د)}] النتيجة سببها صياغة الطلب وحدها، لأن كلمة «عالِم» مذكّرة، ولو كتبت «عالِمة» لجاءت الصور متنوعة الأعراق والعصور تلقائيًا، فالانحياز كامن في اللغة العربية لا في بيانات التدريب.
  \item[\textbf{هـ)}] تحدد في الطلب السياق التاريخي والتنوع المطلوب، ثم تعتمد أول صور مطابقة للوصف دون مراجعة، لأن تحديد السياق في الطلب يضمن الدقة التاريخية في اللباس والأدوات والعمارة ويغني عن أي تدقيق لاحق.
\end{itemize}
\par\vspace{6pt}
\noindent
{\large\khatt\textbf{الوضعية 4}}\hfill{\small(الفصل 5)}\par\nopagebreak\vspace{3pt}
\setlength{\fboxsep}{7pt}\colorbox{gray!10}{\parbox{\dimexpr\linewidth-14pt}{وجد طالب لوفاة شاعر عباسي تاريخين مختلفين: الأول في كتاب تراجم من القرن السابع الهجري، والثاني في موقع تعليمي معاصر لا يذكر مصدره. ثم سأل أداة ذكاء اصطناعي، فقدّمت تاريخًا ثالثًا بثقة ودون إحالة. ومشرفه يشترط توثيق كل معلومة. ما المنهج الأسلم للوصول إلى معلومة يمكن الدفاع عنها في بحثه؟}}\par\nopagebreak
\begin{itemize}[leftmargin=1.0cm, labelsep=6pt, itemsep=3pt, topsep=4pt, parsep=0pt]
  \item[\textbf{أ)}] يعتمد تاريخ الموقع المعاصر، لأن المصادر الحديثة تستفيد مما سبقها فهي أدق من القديمة دائمًا، ويستبعد تاريخ الأداة لأنه بلا مصدر، ولا يذكر الخلاف في بحثه حتى لا يشوّش على القارئ.
  \item[\textbf{ب)}] يعتمد تاريخ الأداة لأنها اطلعت على المصادر القديمة والحديثة معًا ورجّحت بينها، ويذكر في الهامش أنه مأخوذ منها، لأن الإفصاح عن مصدر المعلومة يجعلها موثقة علميًا.
  \item[\textbf{ج)}] يحسب التاريخ الواقع في الوسط بين التواريخ الثلاثة ويعتمده، لأن التوسط بين الروايات يقلّل احتمال الخطأ، ثم يذكر المصادر الثلاثة في الهامش ليبيّن أنه رجع إليها جميعًا قبل الترجيح.
  \item[\textbf{د)}] يرجع إلى مصادر التراجم المتقدمة الأخرى ويقارن رواياتها، ويبحث عن مصدر الموقع المعاصر، ثم يرجّح بالأقدم والأوثق مع ذكر الخلاف إن بقي، ولا يعتمد تاريخ الأداة لأنه بلا إسناد يمكن تتبعه.
  \item[\textbf{هـ)}] يعتمد تاريخ كتاب التراجم القديم مباشرة دون الرجوع إلى مصادر أخرى، لأن أي مصدر قديم أوثق بالضرورة من أي مصدر حديث، ويستبعد الموقع والأداة معًا دون البحث عن مصدر كل منهما، وهي قاعدة معروفة في التحقيق.
\end{itemize}
\par\vspace{6pt}
\noindent
{\large\khatt\textbf{الوضعية 5}}\hfill{\small(الفصل 6)}\par\nopagebreak\vspace{3pt}
\setlength{\fboxsep}{7pt}\colorbox{gray!10}{\parbox{\dimexpr\linewidth-14pt}{أنتجت أستاذة بأداة تحويل النص إلى صوت قراءة مسجلة لقصيدة عمودية مقررة، لتعرضها على تلاميذها نموذجًا للإلقاء. عند الاستماع لاحظت أخطاء في نطق أواخر بعض الكلمات، ووقفات في غير موضعها تكسر الوزن، بينما جاء نطق باقي القصيدة سليمًا. وستعرضه أيضًا على تلاميذ قسمين آخرين. ما الإجراء الأنسب قبل عرض التسجيل؟}}\par\nopagebreak
\begin{itemize}[leftmargin=1.0cm, labelsep=6pt, itemsep=3pt, topsep=4pt, parsep=0pt]
  \item[\textbf{أ)}] تعرض التسجيل كما هو على أنه نموذج للإلقاء، ثم تطلب من التلاميذ اكتشاف أخطائه في حصة لاحقة، لأن سماع النموذج الخاطئ أولًا ثم تصحيحه يرسّخ النطق الصحيح أكثر من سماع نموذج سليم.
  \item[\textbf{ب)}] تقطع المقاطع الخاطئة من التسجيل وتعرض الباقي، لأن حذف الأبيات التي فيها أخطاء أسرع من تصحيحها، ولا يضر التلاميذ أن يسمعوا القصيدة ناقصة ما دام ما بقي منها سليمًا في النطق والوزن.
  \item[\textbf{ج)}] تضبط النص بالشكل التام وتضيف علامات الوقف عند نهايات الأشطر، ثم تعيد التوليد وتراجع التسجيل بيتًا بيتًا، وتسجّل بصوتها المواضع التي بقيت خاطئة، لأن النموذج المقدَّم للتلاميذ يجب أن يكون سليمًا كله.
  \item[\textbf{د)}] تحذف الحركات من النص وتعيد التوليد، لأن الأدوات الصوتية الحديثة تُعرب الكلمات بنفسها حسب السياق، والحركات المكتوبة تربكها، فتنتج قراءة أسلم في أواخر الكلمات وفي الوقفات الوزنية.
  \item[\textbf{هـ)}] تضبط النص بالشكل التام ثم تعيد التوليد، وتعتمد التسجيل الجديد مباشرة، لأن الضبط التام يضمن وحده نطقًا سليمًا لأواخر الكلمات والوقفات الوزنية معًا، فلا حاجة إلى الاستماع إليه مرة أخرى.
\end{itemize}
\par\vspace{6pt}
\begin{center}
{\khatt\textbf{جدول الإجابات}}\\[3pt]
\renewcommand{\arraystretch}{1.6}
\begin{tabular}{|>{\centering\arraybackslash}m{1.9cm}|*{5}{>{\centering\arraybackslash}m{1.2cm}|}}
\hline
\textbf{رقم السؤال} & \textbf{1} & \textbf{2} & \textbf{3} & \textbf{4} & \textbf{5} \\ \hline
\textbf{الإجابة} &  &  &  &  &  \\[8pt] \hline
\end{tabular}

\vspace{8pt}
\renewcommand{\arraystretch}{2.2}
\begin{tabular}{|>{\centering\arraybackslash}m{7.5cm}|>{\centering\arraybackslash}m{3.2cm}|>{\centering\arraybackslash}m{3.2cm}|}
\hline
\textbf{الاسم واللقب} & \textbf{الفوج} & \textbf{العلامة} \\ \hline
 & & \hfill /20 \\ \hline
\end{tabular}
\end{center}
\end{document}
